% =========================================================
% Continuity — Canonical Notes
% Chapter: continuity
% Volume III · Analysis
% =========================================================

\section{Continuity}
\label{sec:continuity}

% ---------------------------------------------------------
% Toolkit
% ---------------------------------------------------------

\begin{tcolorbox}[title={Toolkit — Continuity}]
\begin{itemize}
  \item Three equivalent forms: $\varepsilon$-$\delta$, neighbourhood, sequential.
  \item At a cluster point $c$: continuity $\iff$ $\lim_{x \to c} f(x) = f(c)$.
  \item At an isolated point: continuity holds vacuously.
  \item Algebra of continuous functions: sums, products, scalar multiples, quotients (denominator nonzero), compositions.
  \item $C^{0}(X)$ denotes the vector space of all continuous real-valued functions on $X$.
\end{itemize}
\end{tcolorbox}

% ---------------------------------------------------------
% Definition — Relative Neighborhood
% ---------------------------------------------------------

\begin{tcolorbox}[title={Definition — Relative Neighborhood}]
\begin{definition}[Relative Neighborhood]
\label{def:relative-neighborhood}
Let $E \subseteq \mathbb{R}$ and let $a \in E$. A \textbf{neighborhood
of $a$ in $E$}, denoted $U_{E}(a)$, is a set of the form
\[
  U_{E}(a) \;:=\; E \cap U(a),
\]
where $U(a)$ is a neighborhood of $a$ in $\mathbb{R}$.
\end{definition}
\end{tcolorbox}

\begin{remark*}[Standard quantified statement]
\[
  U_{E}(a) = E \cap (a - \delta,\, a + \delta)
  \quad \text{for some } \delta > 0.
\]
\end{remark*}

\begin{remark*}[Interpretation]
A relative neighborhood is the portion of the domain $E$ that lies
inside a full $\delta$-ball around $a$. The construction restricts
the ambient neighborhood to the domain, enabling local properties
to be stated on arbitrary subsets of $\mathbb{R}$ without
presupposing that a full ball around $a$ lies in $E$.
\end{remark*}

% ---------------------------------------------------------
% Definition — Continuity at a Point (ε-δ form)
% ---------------------------------------------------------

\begin{tcolorbox}[title={Definition — Continuity at a Point}]
\begin{definition}[Continuity at a Point]
\label{def:continuous-at-point}
Let $A \subseteq \mathbb{R}$, let $f : A \to \mathbb{R}$, and let
$c \in A$. The function $f$ is \textbf{continuous at $c$} if for every
$\varepsilon > 0$ there exists $\delta > 0$ such that for every $x \in A$,
\[
  |x - c| < \delta \;\Rightarrow\; |f(x) - f(c)| < \varepsilon.
\]
If $f$ is not continuous at $c$, then $f$ is \textbf{discontinuous at
$c$}.
\end{definition}
\end{tcolorbox}

\begin{remark*}[Standard quantified statement]
\[
  \operatorname{ContinuousAt}(f, c)
  \;\colonequiv\;
  \forall \varepsilon > 0,\;
  \exists \delta > 0,\;
  \forall x \in A :\;
  |x - c| < \delta \;\Rightarrow\; |f(x) - f(c)| < \varepsilon.
\]
\end{remark*}

\begin{remark*}[Negated quantified statement]
\[
  \neg\,\operatorname{ContinuousAt}(f, c)
  \;\colonequiv\;
  \exists \varepsilon > 0,\;
  \forall \delta > 0,\;
  \exists x \in A :\;
  |x - c| < \delta \;\wedge\; |f(x) - f(c)| \geq \varepsilon.
\]
\end{remark*}

\begin{remark*}[Failure mode decomposition]
Discontinuity at $c$ means some fixed output gap $\varepsilon_0 > 0$
survives every shrinking of the input window. Three canonical failure
modes:
\begin{itemize}
  \item \textbf{Jump:} both one-sided limits exist but differ from $f(c)$.
  \item \textbf{Removable:} two-sided limit exists but $\lim_{x \to c} f(x) \neq f(c)$.
  \item \textbf{Essential:} at least one one-sided limit fails to exist.
\end{itemize}
\end{remark*}

\begin{remark*}[Interpretation]
Two structural differences from the limit definition distinguish
continuity. First, $x = c$ is permitted — the domain condition is
$|x - c| < \delta$, not $0 < |x - c| < \delta$. Second, the target
value is $f(c)$ itself, not a separate $L$. At a cluster point $c$,
continuity is therefore exactly the condition $\lim_{x \to c} f(x) = f(c)$.
At an isolated point, every $\delta$-ball meets $A$ only at $c$, so
the implication holds vacuously for every $\varepsilon > 0$.

\smallskip
\noindent
\hyperref[prf:equivalent-forms-continuity]{\textit{Go to proof (Three Equivalent Forms).}}
\end{remark*}

% ---------------------------------------------------------
% Definition — Continuity at a Point (neighbourhood form)
% ---------------------------------------------------------

\begin{definition}[Continuity at a Point — Neighbourhood Form]
\label{def:continuous-at-point-nbhd}
Let $A \subseteq \mathbb{R}$, let $f : A \to \mathbb{R}$, and let
$c \in A$. The function $f$ is \textbf{continuous at $c$} (neighbourhood
form) if for every $\varepsilon$-neighbourhood $V_{\varepsilon}(f(c))$
of $f(c)$ there exists a $\delta$-neighbourhood $V_{\delta}(c)$ of $c$
such that
\[
  f\!\bigl(A \cap V_{\delta}(c)\bigr) \subseteq V_{\varepsilon}(f(c)).
\]
\end{definition}

\begin{remark*}[Standard quantified statement]
\[
  \forall \varepsilon > 0,\;
  \exists \delta > 0 :\;
  \forall x \in A \cap V_{\delta}(c),\;
  f(x) \in V_{\varepsilon}(f(c)).
\]
\end{remark*}

\begin{remark*}[Interpretation]
The neighbourhood form replaces absolute values with set-membership.
The image of the full $\delta$-slice of $A$ must land inside the
$\varepsilon$-ball around $f(c)$. This differs from the limit
neighbourhood form in one place: the full neighbourhood
$V_{\delta}(c)$ replaces the punctured neighbourhood
$V^{*}_{\delta}(c)$, consistent with $x = c$ being permitted.
\end{remark*}

% ---------------------------------------------------------
% Definition — Continuity at a Point (sequential form)
% ---------------------------------------------------------

\begin{definition}[Continuity at a Point — Sequential Form]
\label{def:continuous-at-point-seq}
Let $A \subseteq \mathbb{R}$, let $f : A \to \mathbb{R}$, and let
$c \in A$. The function $f$ is \textbf{continuous at $c$} (sequential
form) if for every sequence $(x_{n}) \subseteq A$ with $x_{n} \to c$,
\[
  f(x_{n}) \to f(c).
\]
\end{definition}

\begin{remark*}[Standard quantified statement]
\[
  \forall (x_{n}) \subseteq A,\;
  x_{n} \to c
  \;\Rightarrow\;
  f(x_{n}) \to f(c).
\]
\end{remark*}

\begin{remark*}[Interpretation]
The sequential form permits $x_{n} = c$, in contrast to the sequential
criterion for function limits, which required $x_{n} \neq c$.
Continuity transports every convergent sequence in $A$ to a convergent
sequence in $\mathbb{R}$; the limit of the transported sequence is the
value of $f$ at the limit of the original.
\end{remark*}

% ---------------------------------------------------------
% Theorem — Three Equivalent Forms of Continuity
% ---------------------------------------------------------

\begin{theorem}[Three Equivalent Forms of Continuity]
\label{thm:continuity-three-equiv}
Let $A \subseteq \mathbb{R}$, let $f : A \to \mathbb{R}$, and let
$c \in A$. The following are equivalent:
\begin{enumerate}[label=\textup{(\roman*)}]
  \item $\forall \varepsilon > 0,\;\exists \delta > 0,\;\forall x \in A:\;
    |x - c| < \delta \Rightarrow |f(x) - f(c)| < \varepsilon$.
    \hfill ($\varepsilon$-$\delta$)
  \item $\forall V_{\varepsilon}(f(c)),\;\exists V_{\delta}(c):\;
    f(A \cap V_{\delta}(c)) \subseteq V_{\varepsilon}(f(c))$.
    \hfill (neighbourhood)
  \item $\forall (x_{n}) \subseteq A,\;
    x_{n} \to c \Rightarrow f(x_{n}) \to f(c)$.
    \hfill (sequential)
\end{enumerate}

\smallskip
\noindent
\hyperref[prf:equivalent-forms-continuity]{\textit{Go to proof.}}
\end{theorem}

\begin{remark*}[Standard quantified statement]
\[
\begin{aligned}
(i)   &\quad \forall \varepsilon > 0,\;
             \exists \delta > 0,\;
             \forall x \in A:\;
             |x - c| < \delta \Rightarrow |f(x) - f(c)| < \varepsilon \\
(ii)  &\quad \forall \varepsilon > 0,\;
             \exists \delta > 0:\;
             f(A \cap V_{\delta}(c)) \subseteq V_{\varepsilon}(f(c)) \\
(iii) &\quad \forall (x_{n}) \subseteq A,\;
             x_{n} \to c \Rightarrow f(x_{n}) \to f(c)
\end{aligned}
\]
\end{remark*}

\begin{remark*}[Interpretation]
$(i) \iff (ii)$ is a notation translation: $|x - c| < \delta$ means
$x \in V_{\delta}(c)$ and $|f(x) - f(c)| < \varepsilon$ means
$f(x) \in V_{\varepsilon}(f(c))$. $(i) \iff (iii)$ mirrors the
sequential criterion for limits: the forward direction feeds $\delta$
to the convergence of $(x_{n})$; the backward direction is contrapositive
with the $1/n$ trick to construct a bad sequence from failure of
$\varepsilon$-$\delta$.
\end{remark*}

% ---------------------------------------------------------
% Definition — Discontinuity Criterion
% ---------------------------------------------------------

\begin{definition}[Discontinuity Criterion]
\label{def:discontinuity-criterion}
Let $A \subseteq \mathbb{R}$, let $f : A \to \mathbb{R}$, and let
$c \in A$. The function $f$ is \textbf{discontinuous at $c$} if and
only if there exists a sequence $(x_{n}) \subseteq A$ such that
$x_{n} \to c$ but $f(x_{n}) \not\to f(c)$.
\end{definition}

\begin{remark*}[Standard quantified statement]
\[
  \neg\,\operatorname{ContinuousAt}(f, c)
  \;\iff\;
  \exists (x_{n}) \subseteq A :\;
  x_{n} \to c
  \;\wedge\;
  f(x_{n}) \not\to f(c).
\]
\end{remark*}

\begin{remark*}[Interpretation]
One witness sequence suffices to establish discontinuity. The sequence
must converge; its image under $f$ must fail to converge to $f(c)$.
Continuity, by contrast, demands that every such sequence have image
converging to $f(c)$ — a universal condition.
\end{remark*}

% ---------------------------------------------------------
% Definition — Continuity on a Set
% ---------------------------------------------------------

\begin{tcolorbox}[title={Definition — Continuity on a Set}]
\begin{definition}[Continuity on a Set]
\label{def:continuity-on-set}
Let $A \subseteq \mathbb{R}$, let $f : A \to \mathbb{R}$, and let
$B \subseteq A$. The function $f$ is \textbf{continuous on $B$} if
$f$ is continuous at every point of $B$.
\end{definition}
\end{tcolorbox}

\begin{remark*}[Standard quantified statement]
\[
  \operatorname{ContinuousOn}(f, B)
  \;\colonequiv\;
  \forall c \in B,\;
  \forall \varepsilon > 0,\;
  \exists \delta > 0,\;
  \forall x \in A :\;
  |x - c| < \delta \Rightarrow |f(x) - f(c)| < \varepsilon.
\]
\end{remark*}

\begin{remark*}[Interpretation]
Continuity on a set is a pointwise condition: $\delta$ may depend on
both $\varepsilon$ and the base point $c$. This is the key distinction
from uniform continuity, where $\delta$ depends only on $\varepsilon$.
\end{remark*}

% ---------------------------------------------------------
% Theorem — Algebra of Continuity at a Point
% ---------------------------------------------------------

\begin{theorem}[Algebra of Continuity at a Point]
\label{thm:continuity-algebra-at-a-point}
Let $A \subseteq \mathbb{R}$, let $f, g : A \to \mathbb{R}$, let
$b \in \mathbb{R}$, and let $c \in A$. Assume $f$ and $g$ are
continuous at $c$.
\begin{enumerate}[label=\textup{(\alph*)}]
  \item The functions $f + g$, $f - g$, $fg$, and $bf$ are continuous
    at $c$.
  \item If $h : A \to \mathbb{R}$ is continuous at $c$ and $h(x) \neq 0$
    for all $x \in A$, then $f/h$ is continuous at $c$.
\end{enumerate}

\smallskip
\noindent
\hyperref[prf:algebra-of-continuity-at-a-point]{\textit{Go to proof.}}
\end{theorem}

\begin{remark*}[Standard quantified statement]
\[
\begin{aligned}
(a)\quad &
\bigl(\operatorname{ContinuousAt}(f,c)
\;\wedge\;
\operatorname{ContinuousAt}(g,c)\bigr)
\;\Rightarrow\;
\operatorname{ContinuousAt}(f+g,\,c)
\;\wedge\;
\operatorname{ContinuousAt}(fg,\,c). \\[4pt]
(b)\quad &
\bigl(\operatorname{ContinuousAt}(f,c)
\;\wedge\;
\operatorname{ContinuousAt}(h,c)
\;\wedge\;
\forall x \in A,\; h(x) \neq 0\bigr)
\;\Rightarrow\;
\operatorname{ContinuousAt}(f/h,\,c).
\end{aligned}
\]
\end{remark*}

\begin{remark*}[Interpretation]
The algebra of continuity reduces to the algebra of limits at a point.
Each operation on continuous functions produces a continuous function,
provided the operation is defined. The denominator condition in (b)
is global on $A$, not merely local at $c$, because $f/h$ must be
defined everywhere on $A$.
\end{remark*}

% ---------------------------------------------------------
% Theorem — Algebra of Continuity on a Set
% ---------------------------------------------------------

\begin{theorem}[Algebra of Continuity on a Set]
\label{thm:continuity-algebra-on-a-set}
Let $A \subseteq \mathbb{R}$, let $f, g : A \to \mathbb{R}$, and let
$b \in \mathbb{R}$. Assume $f$ and $g$ are continuous on $A$.
\begin{enumerate}[label=\textup{(\alph*)}]
  \item The functions $f + g$, $f - g$, $fg$, and $bf$ are continuous
    on $A$.
  \item If $h : A \to \mathbb{R}$ is continuous on $A$ and $h(x) \neq 0$
    for all $x \in A$, then $f/h$ is continuous on $A$.
\end{enumerate}

\smallskip
\noindent
\hyperref[prf:algebra-of-continuity-on-a-set]{\textit{Go to proof.}}
\end{theorem}

\begin{remark*}[Standard quantified statement]
\[
  \operatorname{ContinuousOn}(f,A)
  \;\wedge\;
  \operatorname{ContinuousOn}(g,A)
  \;\Rightarrow\;
  \operatorname{ContinuousOn}(f+g,\,A)
  \;\wedge\;
  \operatorname{ContinuousOn}(fg,\,A).
\]
\end{remark*}

\begin{remark*}[Interpretation]
This is the pointwise algebra theorem applied uniformly at every
$c \in A$. No new argument is needed beyond the at-a-point version.
\end{remark*}

% ---------------------------------------------------------
% Theorem — Composition of Continuous Functions
% ---------------------------------------------------------

\begin{theorem}[Composition of Continuous Functions]
\label{thm:composition-of-continuous-functions}
Let $X, Y \subseteq \mathbb{R}$, let $f : X \to \mathbb{R}$ and
$g : Y \to \mathbb{R}$ with $f(X) \subseteq Y$. If $f$ is continuous
at $x_{0} \in X$ and $g$ is continuous at $f(x_{0}) \in Y$, then
the composition $h = g \circ f : X \to \mathbb{R}$ is continuous at
$x_{0}$.

\smallskip
\noindent
\hyperref[prf:composition-is-continuous]{\textit{Go to proof.}}
\end{theorem}

\begin{remark*}[Standard quantified statement]
\[
  \operatorname{ContinuousAt}(f,\,x_{0})
  \;\wedge\;
  \operatorname{ContinuousAt}(g,\,f(x_{0}))
  \;\Rightarrow\;
  \operatorname{ContinuousAt}(g \circ f,\,x_{0}).
\]
\end{remark*}

\begin{remark*}[Interpretation]
Continuity is preserved under composition. The proof is a two-step
$\varepsilon$-$\delta$ argument: choose $\delta_{1}$ from continuity
of $g$ at $f(x_{0})$ for input $\varepsilon$, then choose $\delta$
from continuity of $f$ at $x_{0}$ for input $\delta_{1}$. The chain
$|x - x_{0}| < \delta \Rightarrow |f(x) - f(x_{0})| < \delta_{1}
\Rightarrow |g(f(x)) - g(f(x_{0}))| < \varepsilon$ closes the
argument.
\end{remark*}

% ---------------------------------------------------------
% Theorem — Local Properties of Continuous Functions
% ---------------------------------------------------------

\begin{theorem}[Local Properties of Continuous Functions]
\label{thm:local-properties-of-continuous-functions}
Let $E \subseteq \mathbb{R}$ and let $f : E \to \mathbb{R}$ be
continuous at $a \in E$. Then:
\begin{enumerate}[label=\arabic*$^{\circ}$]
  \item \textbf{Local boundedness.} $f$ is bounded in some
    neighborhood $U_{E}(a)$.
  \item \textbf{Local sign stability.} If $f(a) \neq 0$, then all
    values of $f$ in some neighborhood $U_{E}(a)$ have the same sign
    as $f(a)$.
  \item \textbf{Algebra at a point.} Sums, differences, and scalar
    multiples of functions continuous at $a$ are continuous at $a$;
    quotients are continuous at $a$ when the denominator is nonzero
    there.
  \item \textbf{Composition.} If $g : Y \to \mathbb{R}$ is continuous
    at $b \in Y$ and $f : E \to Y$ satisfies $f(a) = b$, then
    $g \circ f$ is continuous at $a$.
\end{enumerate}

\smallskip
\noindent
\hyperref[prf:local-properties-of-continuous-functions]{\textit{Go to proof.}}
\end{theorem}

\begin{remark*}[Interpretation]
This theorem packages four separate local results into a single
statement following Zorich. Parts 3 and 4 restate the Algebra and
Composition theorems in local form. Part 2 is Sign Preservation
(\hyperref[prop:sign-preservation]{Proposition}). Part 1 is the
pointwise precursor to the global Boundedness Theorem: continuity
at a single point forces local finiteness of the function values.
\end{remark*}

% ---------------------------------------------------------
% Lemma — Power Function Is Continuous
% ---------------------------------------------------------

\begin{lemma}[Power Function Is Continuous]
\label{lem:power-function-continuous}
Let $n \in \mathbb{N}$. The function $f : \mathbb{R} \to \mathbb{R}$
defined by $f(x) = x^{n}$ is continuous.

\smallskip
\noindent
\hyperref[prf:power-function-is-continuous]{\textit{Go to proof.}}
\end{lemma}

\begin{remark*}[Standard quantified statement]
\[
  \forall n \in \mathbb{N},\;
  \forall x_{0} \in \mathbb{R},\;
  \forall \varepsilon > 0,\;
  \exists \delta > 0,\;
  \forall x \in \mathbb{R} :\;
  |x - x_{0}| < \delta \Rightarrow |x^{n} - x_{0}^{n}| < \varepsilon.
\]
\end{remark*}

\begin{remark*}[Interpretation]
Every monomial is continuous on $\mathbb{R}$. The proof factors
$x^{n} - x_{0}^{n} = (x - x_{0})(x^{n-1} + x^{n-2}x_{0} + \cdots
+ x_{0}^{n-1})$ and estimates the second factor on a bounded
neighborhood of $x_{0}$.
\end{remark*}

% ---------------------------------------------------------
% Theorem — Polynomials Are Continuous
% ---------------------------------------------------------

\begin{theorem}[Any Real Polynomial Is a Continuous Function]
\label{thm:polynomials-are-continuous}
Let $P(x) = a_{0} + a_{1}x + \cdots + a_{n}x^{n}$ be a real
polynomial. Then $P$ is continuous on $\mathbb{R}$.

\smallskip
\noindent
\hyperref[prf:polynomials-are-continuous]{\textit{Go to proof.}}
\end{theorem}

\begin{remark*}[Standard quantified statement]
\[
  \forall n \in \mathbb{N},\;
  \forall a_{0}, \ldots, a_{n} \in \mathbb{R},\;
  \operatorname{ContinuousOn}(P, \mathbb{R}).
\]
\end{remark*}

\begin{remark*}[Interpretation]
Continuity of polynomials follows by finitely many applications of
the Algebra of Continuity theorem: each monomial $x^{k}$ is
continuous, and continuity is preserved under addition and scalar
multiplication.
\end{remark*}

% ---------------------------------------------------------
% Proposition — Rational Functions Are Continuous on Their Domain
% ---------------------------------------------------------

\begin{proposition}[Rational Functions Are Continuous on Their Domain]
\label{prop:rational-functions-continuous}
Let $P$ and $Q$ be real polynomials and let
$D = \{x \in \mathbb{R} : Q(x) \neq 0\}$.
Then the rational function $R(x) = P(x)/Q(x)$ is continuous on $D$.

\smallskip
\noindent
\hyperref[prf:rational-function-continuous-on-domain]{\textit{Go to proof.}}
\end{proposition}

\begin{remark*}[Standard quantified statement]
\[
  \forall x_{0} \in D,\;
  \operatorname{ContinuousAt}(R, x_{0}).
\]
\end{remark*}

\begin{remark*}[Interpretation]
Polynomials are continuous everywhere; the quotient rule for continuous
functions applies at every point where the denominator is nonzero.
The domain $D$ is the natural domain of $R$ — no continuity claim is
made at zeros of $Q$.
\end{remark*}

% ---------------------------------------------------------
% Theorem — Square Root Preserves Continuity
% ---------------------------------------------------------

\begin{theorem}[Square Root Preserves Continuity]
\label{thm:square-root-preserves-continuity}
Let $A \subseteq \mathbb{R}$, let $f : A \to \mathbb{R}$ with
$f(x) \geq 0$ for all $x \in A$, and define $(\sqrt{f})(x) :=
\sqrt{f(x)}$.
\begin{enumerate}[label=\textup{(\alph*)}]
  \item If $f$ is continuous at $c \in A$, then $\sqrt{f}$ is
    continuous at $c$.
  \item If $f$ is continuous on $A$, then $\sqrt{f}$ is continuous
    on $A$.
\end{enumerate}

\smallskip
\noindent
\hyperref[prf:square-root-preserves-continuity]{\textit{Go to proof.}}
\end{theorem}

\begin{remark*}[Interpretation]
The square-root operation preserves continuity on nonnegative-valued
functions. Part (b) follows immediately by applying part (a) at each
point of $A$.
\end{remark*}

% ---------------------------------------------------------
% Proposition — Exponentiation Is Continuous, I
% ---------------------------------------------------------

\begin{proposition}[Exponentiation Is Continuous, I]
\label{prop:exponentiation-continuous-base}
Let $a > 0$. The function $f : \mathbb{R} \to \mathbb{R}$ defined by
$f(x) := a^{x}$ is continuous.

\smallskip
\noindent
\hyperref[prf:exponentiation-is-continuous-1]{\textit{Go to proof.}}
\end{proposition}

\begin{remark*}[Standard quantified statement]
\[
  \forall a > 0,\;
  \operatorname{ContinuousOn}(x \mapsto a^{x},\, \mathbb{R}).
\]
\end{remark*}

\begin{remark*}[Interpretation]
For a fixed positive base, exponentiation in the exponent is continuous
on all of $\mathbb{R}$.
\end{remark*}

% ---------------------------------------------------------
% Proposition — Exponentiation Is Continuous, II
% ---------------------------------------------------------

\begin{proposition}[Exponentiation Is Continuous, II]
\label{prop:exponentiation-continuous-exponent}
Let $p \in \mathbb{R}$. The function $f : (0, \infty) \to \mathbb{R}$
defined by $f(x) := x^{p}$ is continuous.
\end{proposition}

\begin{remark*}[Standard quantified statement]
\[
  \forall p \in \mathbb{R},\;
  \operatorname{ContinuousOn}(x \mapsto x^{p},\, (0,\infty)).
\]
\end{remark*}

\begin{remark*}[Interpretation]
For a fixed real exponent, the power function in the base is
continuous on the positive reals.
\end{remark*}

% ---------------------------------------------------------
% Proposition — Sign Preservation
% ---------------------------------------------------------

\begin{proposition}[Sign Preservation]
\label{prop:sign-preservation}
Let $f : E \to \mathbb{R}$ be continuous at $x_{0} \in E$.
\begin{itemize}
  \item If $f(x_{0}) > 0$, there exists $\delta > 0$ such that
    $f(x) > 0$ for all $x \in E$ with $|x - x_{0}| < \delta$.
  \item If $f(x_{0}) < 0$, there exists $\delta > 0$ such that
    $f(x) < 0$ for all $x \in E$ with $|x - x_{0}| < \delta$.
\end{itemize}

\smallskip
\noindent
\hyperref[prf:sign-preservation]{\textit{Go to proof.}}
\end{proposition}

\begin{remark*}[Standard quantified statement]
\[
  \operatorname{ContinuousAt}(f, x_{0})
  \;\wedge\;
  f(x_{0}) > 0
  \;\Rightarrow\;
  \exists \delta > 0,\;
  \forall x \in E :\;
  |x - x_{0}| < \delta \Rightarrow f(x) > 0.
\]
\end{remark*}

\begin{remark*}[Interpretation]
Continuity at a point with nonzero value forces sign stability on a
neighborhood. The standard proof applies the $\varepsilon$-$\delta$
definition with $\varepsilon = |f(x_{0})|/2$; the resulting
$\delta$-neighborhood keeps $f$ bounded away from zero on the same
side as $f(x_{0})$.
\end{remark*}

% ---------------------------------------------------------
% Proposition — Continuous Extension to a Closed Interval
% ---------------------------------------------------------

\begin{proposition}[Continuous Extension to a Closed Interval]
\label{prop:continuous-extension-closed-interval}
Let $f : (a, b] \to \mathbb{R}$ be continuous and suppose
$\lim_{x \to a^{+}} f(x)$ exists. Define
\[
  \widetilde{f}(x) :=
  \begin{cases}
    \displaystyle\lim_{x \to a^{+}} f(x), & x = a, \\[4pt]
    f(x), & x \in (a, b].
  \end{cases}
\]
Then $\widetilde{f} : [a, b] \to \mathbb{R}$ is continuous.

\smallskip
\noindent
\hyperref[prf:continuous-extension-to-closed-interval]{\textit{Go to proof.}}
\end{proposition>

\begin{remark*}[Interpretation]
If a one-sided limit exists at the missing endpoint, the function can
be continuously filled in there. Continuity away from $a$ is inherited
from $f$; continuity at $a$ uses the definition of $\widetilde{f}(a)$
together with one-sided continuity.
\end{remark*}

% ---------------------------------------------------------
% Definition — Space of Continuous Functions
% ---------------------------------------------------------

\begin{definition}[Space of Continuous Functions]
\label{def:space-of-continuous-functions}
Let $X \subseteq \mathbb{R}$. The \textbf{space of continuous
real-valued functions on $X$} is
\[
  C^{0}(X;\mathbb{R})
  \;:=\;
  \{\, f : X \to \mathbb{R} : f \text{ is continuous on } X \,\}.
\]
When the codomain is clear, this space is denoted $C^{0}(X)$.
\end{definition}

\begin{remark*}[Interpretation]
The set $C^{0}(X;\mathbb{R})$ is closed under addition and scalar
multiplication, forming a real vector space. It is the natural
algebraic setting for studying global properties of continuous
functions and for defining function norms such as the uniform norm
$\|f\|_{\infty} = \sup_{x \in X} |f(x)|$.
\end{remark*}
